\documentclass{article}
\usepackage[T1]{fontenc}
\usepackage{lmodern}
\usepackage{microtype}
\setlength{\emergencystretch}{1em}
\usepackage[polish]{babel}
\usepackage{amssymb}
\usepackage{amsmath}
\usepackage[margin=2.5cm]{geometry}
\usepackage{enumitem}
\setlist{itemsep=2pt, topsep=4pt}

\title{Prognoza pogody perceptronem wielowarstwowym\\ \large Konspekt projektu}
\author{Hubert Błonowski \and Jan Zaborski}
\date{}

\newlist{badania}{enumerate}{1}
\setlist[badania]{label=\thesubsection.\arabic*., leftmargin=*, align=left}
\newlist{punkty}{enumerate}{1}
\setlist[punkty]{label=\thesection.\arabic*., leftmargin=*, align=left}

\begin{document}

\maketitle

\section*{Zadanie i dane}
Korzystamy ze zbioru Historical Hourly Weather Data (Kaggle): pomiary godzinowe z 36 miast
(30 z USA i Kanady, 6 z Izraela) od października 2012 do listopada 2017.
Na podstawie dni $D_{t-3}, D_{t-2}, D_{t-1}$ przewidujemy dzień $D_{t+1}$, pomijając dzień $D_t$ (okno \texttt{I I I X O}).
Sieć ma dwa wyjścia:
\begin{itemize}
  \item średnią temperaturę dnia $D_{t+1}$ (trafienie, gdy błąd $\le 2\,^\circ$C),
  \item wystąpienie silnego wiatru, czyli prędkości $\ge 6$~m/s (ocena przez AUC).
\end{itemize}

Dane dzielimy chronologicznie: uczenie do końca 2015 roku, walidacja 2016, test 2017.
Między zbiorami zostawiamy przerwę długości okna, żeby żadne okno nie obejmowało dni z dwóch zbiorów.
Statystyki normalizacji liczymy wyłącznie na zbiorze uczącym.
Model jest perceptronem wielowarstwowym rozwijanym na bazie biblioteki z projektu~1;
przepływ sygnału składamy ręcznie z warstw liniowych i funkcji aktywacji.

\section{Badania}

\subsection{Architektura}
\begin{badania}
  \item Liczba warstw ukrytych i neuronów w warstwach.
  \item Funkcje aktywacji: ReLU, tanh, sigmoid.
  \item Różne aktywacje w kolejnych warstwach a jedna funkcja dla całej sieci.
  \item Warianty podwójnego wyjścia:
    \begin{itemize}
      \item wspólny tułów i dwie głowy (wspólne warstwy ukryte),
      \item dwie osobne sieci (bez współdzielenia wag),
      \item częściowe współdzielenie (wspólna tylko pierwsza warstwa),
      \item waga $\lambda$ w koszcie $L = L_\text{temp} + \lambda L_\text{wiatr}$ (balans dwóch zadań).
    \end{itemize}
  \item Kodowanie wyjść:
    \begin{itemize}
      \item temperatura jako wartość standaryzowana lub zmiana względem $D_{t-1}$,
      \item wiatr jako prawdopodobieństwo (sigmoid) lub regresja prędkości z progiem 6~m/s.
    \end{itemize}
\end{badania}

\subsection{Uczenie}
\begin{badania}
  \item Optymalizatory: SGD, SGD z momentum, Adam.
  \item Funkcje kosztu: MSE, MAE i Huber dla temperatury, BCE dla wiatru.
    Ocena modeli osobnymi metrykami:
    \begin{itemize}
      \item temperatura: odsetek trafień ($|\text{błąd}| \le 2\,^\circ$C), MAE, RMSE,
      \item wiatr: AUC, dodatkowo czułość i precyzja przy wybranym progu.
    \end{itemize}
\end{badania}

\subsection{Inicjalizacja}
\begin{badania}
  \item Inicjalizacja He i Xavier a rozkład jednostajny o stałym zakresie.
\end{badania}

\subsection{Dane}
\begin{badania}
  \item Kodowanie wejść:
    \begin{itemize}
      \item agregaty dobowe (średnia, minimum, maksimum) lub bardziej szczegółowe (dzień i noc osobno),
      \item normalizacja globalna lub osobno dla każdego miasta, standaryzacja,
      \item miasto: one-hot lub położenie na sferze jednostkowej
        $(\cos\varphi\cos\lambda,\ \cos\varphi\sin\lambda,\ \sin\varphi)$,
      \item dzień roku jako $(\sin, \cos)$; godzina jako $(\sin, \cos)$ tylko w wariantach bez agregacji dobowej,
      \item opis pogody: one-hot grup lub udział godzin z danym typem pogody w ciągu dnia,
      \item kierunek wiatru jako $(\sin, \cos)$.
    \end{itemize}
  \item Agregacja okna: spłaszczenie trzech dni, statystyki okna (średnia, trend),
    bloki dni o wspólnych wagach (ten sam koder dla każdego dnia, wyniki łączone).
  \item Braki danych: interpolacja krótkich luk, uzupełnianie średnią, odrzucanie okien z brakami.
  \item Wybór cech przez ablację: usuwamy po jednej grupie cech i mierzymy spadek jakości na zbiorze walidacyjnym.
\end{badania}

\subsection{Pozostałe}
\begin{badania}
  \item Predykcja pośrednia dnia $D_t$: pierwsza sieć przewiduje $D_t$, druga przewiduje $D_{t+1}$
    na podstawie okna i przewidzianego $D_t$. Drugą sieć uczymy na przewidzianym, a nie prawdziwym $D_t$.
  \item Dobór hiperparametrów: badania 1.1--1.4 zmieniają jeden czynnik względem konfiguracji bazowej;
    konfigurację końcową wybieramy losowym przeszukiwaniem (random search) najważniejszych parametrów na zbiorze walidacyjnym.
    Zbiór testowy oceniamy raz, na końcu.
  \item Porównanie z metodami klasycznymi:
    \begin{itemize}
      \item persystencja (prognoza równa ostatniemu znanemu dniowi),
      \item klimatologia (średnia dla danego dnia roku i miasta),
      \item regresja liniowa dla temperatury i logistyczna dla wiatru na tych samych cechach.
    \end{itemize}
\end{badania}

\section{Implementacja}
\begin{punkty}
  \item Skrypty danych w Pythonie: wczytanie, agregaty, braki danych, kodowania, okna, podział chronologiczny.
  \item Funkcja aktywacji wybierana osobno dla każdej warstwy.
  \item Podwójna warstwa wyjściowa: dwie głowy i suma kosztów z wagą $\lambda$.
  \item Optymalizatory: SGD z momentum i Adam.
  \item Nowe funkcje kosztu: MAE, Huber, BCE z opcjonalną wagą klasy (silny wiatr może być rzadki).
  \item Metryki oceny: odsetek trafień, MAE, RMSE, AUC.
  \item Inicjalizacja He i Xavier.
  \item Wczesne zatrzymanie na podstawie kosztu walidacyjnego.
  \item Bloki dni o wspólnych wagach (gradient wspólnego kodera jest sumą gradientów z trzech dni).
  \item Predykcja pośrednia (dwie sieci uczone kolejno).
  \item Zewnętrzny skrypt losowego przeszukiwania hiperparametrów.
  \item Rozszerzenie zapisu konfiguracji z projektu~1 o nowe parametry (wersja, ziarno, podział, parametry).
  \item Opcjonalnie: SIMD w obliczeniach macierzowych, biblioteka Rust wywoływana z Pythona.
\end{punkty}

\section{Hipotezy}
\begin{punkty}
  \item Najlepiej sprawdzą się 2 warstwy ukryte (na podstawie doświadczeń z projektu~1).
  \item Wystarczy około 64 neuronów w warstwie.
  \item Różne aktywacje w kolejnych warstwach nie poprawią wyników względem jednej funkcji.
  \item Adam zbiega szybciej niż SGD, ale końcowa jakość jest podobna.
  \item Inicjalizacje He i Xavier dają lepsze i stabilniejsze uczenie niż rozkład jednostajny.
  \item Huber i MAE dają wyższy odsetek trafień ($\le 2\,^\circ$C), MSE niższe RMSE.
  \item Kodowanie wejść:
    \begin{itemize}
      \item normalizacja osobno dla każdego miasta jest lepsza od globalnej,
      \item położenie na sferze działa co najmniej tak dobrze jak one-hot miasta przy mniejszej liczbie wejść,
      \item minimum, maksimum i trend dają lepsze wyniki niż same średnie,
      \item cykliczne kodowanie dnia roku poprawia przede wszystkim prognozę temperatury.
    \end{itemize}
  \item Bloki dni o wspólnych wagach są co najmniej tak dobre jak spłaszczenie okna przy mniejszej liczbie parametrów.
  \item Ablacja pokaże, że najważniejsze są temperatury z poprzednich dni i dzień roku, a opis pogody i kierunek wiatru wnoszą niewiele.
  \item Predykcja pośrednia dnia $D_t$ nie poprawia prognozy $D_{t+1}$.
  \item MLP jest lepszy od persystencji i klimatologii oraz nieznacznie lepszy od regresji liniowej i logistycznej.
\end{punkty}

\end{document}
